% ==================== 第七章 ====================
\chapter{行动纲领：可操作的政策建议清单}

2024年的一个周末，一位参与国家科技政策制定的学者，在办公室里翻看着最近两年各部门发布的AI相关政策文件。桌上堆积了厚厚一摞——《新一代人工智能发展规划》《人工智能创新发展试验区建设方案》《关于加快场景创新的通知》《生成式人工智能服务管理暂行办法》……

文件很多，表态很积极，但他隐约觉得缺了点什么。

"方向都对，但太抽象了，"他对同事说，"比如'加大投入'，加多少？谁出钱？'重视人才'，具体怎么做？谁负责？"

这正是本章要解决的问题。

\textbf{本章不再分析问题，而是直接给出解决方案清单——一份可以直接进入政策文件的行动纲领。}

前六章已经完成了"为什么"、"是什么"、"怎么办"的论证。本章将所有建议汇总为可直接执行的政策清单，每条建议都明确：具体的行动事项、明确的责任主体、清晰的时间节点、可量化的评估标准。

\textbf{本章定位}：这是针对前六章核心内容的"阶段性政策总结"。后续第八至十二章将分别深入探讨科研赋能、人才安全、投资战略、技术前沿、非传统安全等专题，每章末尾会有针对该领域的专项建议。本章的建议与后续章节互为补充，共同构成完整的政策建议体系。

\textbf{原则}：含糊其辞是政策的坟墓。每条建议都必须具体到可以追责。

\section{战略定位确认}

在列出具体行动项之前，必须首先确认战略定位。因为定位决定力度，力度决定成败。

\textbf{大模型是21世纪"不响的原子弹"——国家竞争的决定性力量。}

这不是修辞夸张。让我们比较一下核武器和大模型：

核武器的威力在于"不用"——它的价值在于威慑，在于让对手不敢轻举妄动。而大模型的威力在于"每天都在用"——它无声无息地渗透进经济、社会、文化的每一个角落，日积月累地改变力量对比。

核武器差距可以通过"最小威慑"弥补——即使只有少量核武器，也足以让对手忌惮。但大模型差距会随时间指数级扩大——因为AI能力会自我增强，强者愈强。

核战争有人阻止——因为后果太可怕，所有人都会努力避免。但认知战争无人察觉——它不像战争，更像是一场缓慢但不可逆的侵蚀。

\textbf{为什么必须像"两弹一星"一样对待大模型？}

这关乎国家生存和民族复兴，不是普通的产业问题。没有任何商业力量能够独自完成，必须举国体制。时间窗口极其有限，错过就是永远落后。差距一旦形成，将体现在国家竞争的每一个维度——经济、军事、科技、文化，无一幸免。

\textbf{结论}：必须以"两弹一星"的战略高度、投入力度、紧迫程度对待大模型发展。

\section{顶层设计行动项}

\textbf{行动项1：成立国家级领导小组}

\textbf{优先级}：P0（最高）

\textbf{责任主体}：中央

\textbf{时间节点}：2025年Q2前完成

\textbf{具体内容}：成立由政治局常委分管的"国家认知基础设施建设领导小组"；成员包括科技部、工信部、发改委、教育部、财政部、国安委、网信办主要负责人；设立专职秘书处，配备50-100名专业人员；建立季度汇报和年度评估制度。

\textbf{评估标准}：领导小组是否正式成立并开始运作？第一次全体会议是否召开？《国家认知基础设施建设规划》是否启动编制？

\textbf{行动项2：发布国家认知基础设施建设规划}

\textbf{优先级}：P0

\textbf{责任主体}：国家发改委牵头，科技部、工信部配合

\textbf{时间节点}：2025年Q3前发布

\textbf{具体内容}：明确2025-2035年建设目标和里程碑；确定投资规模（建议2-3万亿元）；明确资金来源和分配机制；建立项目管理和考核机制。

\textbf{评估标准}：规划是否正式发布？是否获得足额预算支持？配套实施方案是否制定？

\section{算力基础设施行动项}

\textbf{行动项3：国家智算中心建设}

\textbf{优先级}：P0

\textbf{责任主体}：国家发改委牵头，工信部配合，地方政府执行

\textbf{时间节点}：2025年Q4首批3个万卡级中心开工；2027年Q2首个万卡级中心投入运营；2030年底10个万卡级中心全部运营

\textbf{投资规模}：1.2-1.5万亿元（十年累计）

\textbf{布局方案}：西部训练集群布局在内蒙古（2-3个）、甘肃/青海（1-2个）、四川/贵州（2-3个），利用清洁能源优势；东部推理集群布局在北京/天津（2个）、上海/杭州（2个）、深圳/广州（1-2个），贴近应用场景；另设专用安全集群，位置保密，用于敏感任务。

\textbf{评估标准}：开工数量和进度是否达标？算力总量是否按计划增长？能耗和运营效率是否达标？

\textbf{行动项4：算力网络互联}

\textbf{优先级}：P1

\textbf{责任主体}：工信部牵头，三大运营商执行

\textbf{时间节点}：2027年底前完成主干网建设

\textbf{投资规模}：1000-1500亿元

\textbf{具体内容}：
第一，骨干网：400G-1T带宽，连接各智算中心。第二，城域网：100G-400G带宽，连接用户。第三，RDMA网络：支持分布式训练。第四，智能调度系统：跨中心任务调度。

\textbf{评估标准}：
网络覆盖范围和带宽；跨中心训练的延迟和效率；算力利用率。

\section{芯片攻关行动项}

\textbf{优先级}：P0

\textbf{责任主体}：工信部牵头，科技部配合，华为等企业执行

\textbf{时间节点}：
2025-2026：存量优化，软件生态初步形成；2026-2028：国产替代，80\%训练任务使用国产芯片；2028-2035：自主领先，关键指标达到国际先进。

\textbf{投资规模}：5000-6000亿元

\textbf{重点任务}：
第一，芯片设计能力提升（华为Ascend系列）。第二，先进封装技术突破（弥补制程差距）。第三，软件生态建设（类CUDA框架）。第四，算子库完善和优化。第五，开发者社区培育。

\textbf{评估标准}：
国产芯片在智算中心的部署比例；软件迁移的完成度和性能；开发者数量和活跃度。

\section{模型研发行动项}

\textbf{优先级}：P0

\textbf{责任主体}：科技部牵头，新成立国家实验室或指定头部企业执行

\textbf{时间节点}：
2026年Q2：工程正式启动；2027年Q4：V1.0发布，达到GPT-5同等水平；2030年：综合能力世界一流。

\textbf{投资规模}：1000-1500亿元

\textbf{组织模式}：
第一，牵头单位：新设国家实验室或指定企业（如深度求索）。第二，参与单位：阿里、百度、华为、腾讯等头部企业，清华、北大等顶尖高校。第三，资源保障：国家智算中心算力优先保障。第四，知识产权：国家持有核心IP，参与单位享有使用权。

\textbf{评估标准}：
模型能力在国际基准测试上的排名；关键能力（推理、代码、数学）的突破情况；国产化程度（数据、算力、框架）。

\textbf{优先级}：P1

\textbf{责任主体}：各行业主管部门牵头

\textbf{时间节点}：2027年底前完成重点领域模型开发

\textbf{投资规模}：500-800亿元

\textbf{重点领域}：
\textbf{第一}，科研模型（科技部）：辅助科学研究。\textbf{第二}，教育模型（教育部）：个性化学习。\textbf{第三}，医疗模型（卫健委）：辅助诊断、药物研发。\textbf{第四}，政务模型（国办）：政务服务、决策支持。\textbf{第五}，法律模型（司法部）：法律咨询、合同审查。\textbf{第六}，金融模型（金融监管总局）：风控、投研。

\textbf{评估标准}：
各领域模型的部署和应用情况；应用效果评估（效率提升、准确率等）；用户反馈和满意度。

\section{人才培养行动项}

\textbf{优先级}：P0

\textbf{责任主体}：教育部牵头，科技部、人社部配合

\textbf{时间节点}：2025年启动，2030年形成完整体系

\textbf{培养目标}：
对齐研究人才：100人；红队测试人才：1000人；AI安全工程师：10000人；AI治理研究人才：500人；复合型战略人才：100人。

\textbf{具体措施}：
第一，设立"AI安全"本科专业方向（10所高校）。第二，设立"AI治理"硕士项目（5所高校）。第三，设立"对齐研究英才班"（博士层次）。第四，建立职业培训和认证体系。第五，举办"国家AI安全挑战赛"。

\textbf{评估标准}：
各类人才培养数量是否达标；毕业生就业和发展情况；国际排名和影响力。

\textbf{优先级}：P0

\textbf{责任主体}：科技部牵头，各省市人才办配合

\textbf{时间节点}：立即启动，持续推进

\textbf{引进目标}：
引进20-30名在顶级AI实验室有对齐研究经验的华人；引进50-100名资深AI安全工程师；引进20-30名AI治理国际专家。

\textbf{待遇标准}：
顶尖对齐研究者：年薪300-800万元+安家费+科研经费；资深安全工程师：年薪150-300万元；治理专家：年薪100-200万元。

\textbf{配套措施}：
第一，解决住房、子女教育、配偶就业。第二，简化签证和居留手续。第三，提供充足科研启动经费。第四，建立灵活的考核机制。

\textbf{评估标准}：
引进人才数量和质量；引进人才的留任率；引进人才的成果产出。

\section{安全防护行动项}

\textbf{优先级}：P0

\textbf{责任主体}：国家网信办牵头，公安部、国安部配合

\textbf{时间节点}：
2025年底：完成架构设计和试点部署；2026年底：覆盖所有政务和关键基础设施AI应用；2027年底：覆盖所有高风险商业应用。

\textbf{四层架构}：
第一，第一层：系统提示硬化（防止用户覆盖）。第二，第二层：输入净化（检测注入和越狱）。第三，第三层：工具/资源隔离（最小权限）。第四，第四层：输出审计与溯源（风险检测+水印）。

\textbf{评估标准}：
覆盖率是否达标；拒止成功率 $\geq$ 95\%；误杀率 $\leq$ 5\%。

\textbf{优先级}：P1

\textbf{责任主体}：科技部牵头，设立专门机构

\textbf{时间节点}：2026年底前建成并运营

\textbf{具体内容}：
第一，成立"国家AI安全测评中心"。第二，建立大模型安全测评标准。第三，组建专业红队测试队伍。第四，建立漏洞库和情报共享机制。第五，定期发布安全评估报告。

\textbf{测评范围}：
所有在境内提供服务的大模型；政务和关键基础设施使用的AI系统；高风险应用场景。

\textbf{评估标准}：
测评能力是否建立；测评覆盖率；漏洞发现和修复情况。

\section{治理体系行动项}

\textbf{优先级}：P1

\textbf{责任主体}：全国人大法工委牵头，科技部、网信办配合

\textbf{时间节点}：
2025年：立项并启动调研；2026年：形成草案并征求意见；2027年：提请审议并出台。

\textbf{核心内容}：
\textbf{第一}，AI发展与治理的基本原则。\textbf{第二}，监管体制和职责分工。\textbf{第三}，高风险AI系统的管理要求。\textbf{第四}，AI安全的基本制度。\textbf{第五}，数据和算法治理。\textbf{第六}，法律责任和救济机制。

\textbf{评估标准}：
立法进度是否按计划推进 法律出台后的实施效果

\textbf{优先级}：P1

\textbf{责任主体}：中央编办牵头

\textbf{时间节点}：2026年底前完成

\textbf{具体架构}：
首先，一委——国家AI发展与安全委员会（中央层面）。其次，一办——国家AI治理办公室（正部级或副部级）。第三，多中心——国家AI安全测评中心；国家AI标准研究中心；国家AI伦理研究中心；国家AI应急响应中心。

\textbf{评估标准}：
机构是否按计划设立；职能划分是否清晰；协调机制是否有效运作。

\section{国际合作行动项}

\textbf{优先级}：P1

\textbf{责任主体}：外交部牵头，科技部配合

\textbf{时间节点}：争取2025年底前启动

\textbf{对话议题}：
第一，AI安全对齐研究合作。第二，防止AI用于大规模杀伤性武器。第三，核指挥控制系统的AI隔离。第四，AI军备控制框架探讨。第五，学术交流机制化。

\textbf{评估标准}：
对话机制是否建立；对话频率和层级；是否达成任何共识。

\textbf{优先级}：P2

\textbf{责任主体}：国家标准委牵头，科技部、工信部配合

\textbf{时间节点}：持续推进

\textbf{重点领域}：
第一，ISO/IEC JTC 1/SC 42（AI国际标准）。第二，IEEE AI标准工作组。第三，联合国AI治理框架。第四，区域合作标准（如RCEP框架）。

\textbf{评估标准}：
中国专家在国际标准组织中的职位数；中国主导或参与制定的标准数量；国际标准采纳中国方案的比例。

\section{总体时间表}

\begin{table}[H]
\centering
\caption{行动纲领总体时间表}
\small
\begin{tabular}{p{2cm}p{10cm}}
\toprule
\textbf{时间} & \textbf{关键里程碑} \\
\midrule
2025Q2 & 国家认知基础设施建设领导小组成立 \\
2025Q3 & 《国家认知基础设施建设规划》发布 \\
2025Q4 & 首批3个国家智算中心开工；AI安全人才计划启动 \\
2025Q4 & 中美AI安全对话机制启动（争取） \\
2026Q2 & 国家基础大模型工程启动 \\
2026Q4 & AI治理"一委一办多中心"架构建成 \\
2026Q4 & 四层安全网关架构完成政务系统部署 \\
2027Q2 & 首个万卡级国家智算中心投入运营 \\
2027Q4 & 国家基础大模型V1.0发布；《人工智能法》出台 \\
2028Q4 & 5个万卡级智算中心运营；国产芯片占比超80\% \\
2030Q4 & 10+万卡级智算中心；基础模型达世界一流 \\
2035 & 全面建成自主可控的国家认知基础设施 \\
\bottomrule
\end{tabular}
\end{table}

\section{资源需求汇总}

\begin{table}[H]
\centering
\caption{十年投资规模汇总（2025-2035）}
\small
\begin{tabular}{p{4cm}p{3cm}p{5cm}}
\toprule
\textbf{领域} & \textbf{投资规模} & \textbf{主要用途} \\
\midrule
算力基础设施 & 1.2-1.5万亿元 & 智算中心、算力网络、能源配套 \\
芯片攻关 & 5000-6000亿元 & 先进制程、设备材料、软件生态 \\
模型研发 & 2000-3000亿元 & 基础模型、垂直模型、开源生态 \\
人才培养 & 1000-1500亿元 & 高校建设、人才引进、职业培训 \\
应用推广 & 1000-1500亿元 & 政务、教育、医疗等应用 \\
安全体系 & 300-500亿元 & 测评中心、安全网关、应急响应 \\
\midrule
\textbf{总计} & \textbf{2-3万亿元} & \\
\bottomrule
\end{tabular}
\end{table}

\section{风险与应急}

\textbf{主要风险}：
第一，芯片断供升级：美国可能进一步收紧管制。第二，技术路线变化：AI技术可能出现颠覆性变革。第三，人才流失：顶尖人才持续外流。第四，安全事件：重大AI安全事件发生。

\textbf{应急准备}：
第一，保持6-12个月关键芯片战略储备。第二，建立技术路线多元化机制。第三，完善人才留用和回流政策。第四，建立AI安全应急响应机制。

\section{后续章节专项建议索引}

\textbf{本章汇总了前六章的核心政策建议。后续第八至十二章将对专项领域进行深入分析，并在各章末尾提出针对性的政策建议。以下是后续章节专项建议的索引，供读者查阅：}

\begin{table}[H]
\centering
\caption{后续章节专项建议索引}
\small
\begin{tabular}{p{1.5cm}p{3.5cm}p{7.5cm}}
\toprule
\textbf{章节} & \textbf{主题} & \textbf{核心建议要点} \\
\midrule
第八章 & 科研奇点 & AI for Science国家战略；国家科研AI平台建设；自动化实验室试点；跨学科人才培养 \\
\midrule
第九章 & 人才强基 & "人才安全"新维度；人才资产负债表评估；"磁吸-锚定-生态"三位一体战略；人才期权池机制 \\
\midrule
第十章 & 资本赋能 & "认知资本"概念；AI国家主权基金设计；政府投资"三环模型"；举国体制风险防范 \\
\midrule
第十一章 & 技术前沿 & "认知能力阶梯"理论；推理模型发展策略；AI Agent布局；具身智能与世界模型前瞻 \\
\midrule
第十二章 & 安全新疆域 & AI系统性风险框架；金融AI安全"三道防线"；关键基础设施"AI韧性"评估；AGI预警机制 \\
\bottomrule
\end{tabular}
\end{table}

\textbf{建议阅读方式}：本书针对不同读者群体设计了多元化的阅读路径。决策者可重点阅读各章的政策建议部分，参考本索引快速定位关键议题；研究者可深入阅读技术分析和理论框架部分，把握前沿动态与学术脉络；执行层面则可重点关注具体的实施路径和时间节点，将战略规划转化为行动方案。

\section{结语}

\textbf{致决策者}

历史经验表明，在技术革命的关键节点，战略判断和决策果断至关重要。

\textbf{蒸汽机时代}，英国的领先奠定了一个世纪的霸权。\\
\textbf{电气化时代}，美国和德国的崛起改变了世界格局。\\
\textbf{信息革命时代}，硅谷的创新塑造了数字经济版图。\\
\textbf{核武器时代}，"两弹一星"确保了中国的大国地位。

\textbf{大模型革命，我们不能错过。}

这不是一项普通的技术，而是关乎国家认知能力的基础设施。\\
这不是一场普通的竞争，而是决定21世纪国家地位的战略角逐。

\textbf{窗口期有限，行动刻不容缓。}

本书给出了详细的路线图、时间表和资源需求。\\
剩下的，是决策和执行。

\textit{愿这本书能为中国的认知基础设施建设贡献绵薄之力。}\\
\textit{愿我们不负这个时代赋予的机遇与使命。}

